\%============================================================
\chapter{Bài Học: Nhân Quả --- Những Gì Bạn Gieo, Bạn Sẽ Gặt (Cause and Effect --- You Reap What You Sow)}
\begin{center}
  \textit{\color{truyengold}``Every seed of goodness sown today is a flower that will bloom tomorrow.''}\\[4pt]
  \textit{(Mỗi hạt giống thiện lành được gieo hôm nay là một đoá hoa sẽ nở mai sau.)}\\[4pt]
  \small\color{truyenblue}--- Cổ Kim Đông Tây ---
\end{center}
\ngancach

\section{Fleming Và Churchill --- Nhân Quả Xuyên Thế Hệ}
\begin{truyen}{Fleming Và Churchill --- Nhân Quả Xuyên Thế Hệ}{Alexander Fleming và Winston Churchill --- Anh, 1908}
\chuhoa{M}{ột} ngày mưa năm 1908, một nông dân nghèo tên Alec Fleming nghe thấy tiếng kêu cứu từ vũng lầy.\\[4pt]
\textit{(On a rainy day in 1908, a poor farmer named Alec Fleming heard cries for help from a muddy bog.)}

Ông lao xuống cứu một cậu bé đang sắp chết đuối --- không hỏi tên, không mong trả ơn.\\[4pt]
\textit{(He plunged in to save a young boy who was nearly drowning --- asking no name, expecting no reward.)}

Hôm sau, một quý tộc giàu có đến cảm ơn và hỏi: ``Tôi có thể giúp gì cho ông không?'' --- đó là cha của cậu bé.\\[4pt]
\textit{(The next day, a wealthy lord came to thank him and asked: `Can I do anything for you?' --- he was the boy's father.)}

Fleming từ chối nhận tiền nhưng nhận học bổng cho con trai ông --- người con đó trở thành Alexander Fleming, người phát minh ra penicillin.\\[4pt]
\textit{(Fleming refused money but accepted a scholarship for his son --- that son became Alexander Fleming, who discovered penicillin.)}

Cậu bé được cứu năm xưa, lớn lên từ học bổng cũng được tặng, là Winston Churchill --- chính penicillin đã cứu mạng ông sau này.\\[4pt]
\textit{(The boy saved that day, later also given a scholarship by that lord, was Winston Churchill --- penicillin itself would later save his life.)}

\end{truyen}
\begin{baihoc}
  Một hành động tốt gieo ra thường trở về nhiều hơn bạn có thể hình dung --- và thường là khi bạn quên nó từ lâu rồi.\\[4pt]
  \textit{(A good deed sown often returns far more than you can imagine --- and usually long after you have forgotten it.)}
\end{baihoc}

\section{Cụ Già Và Cây Hồ Đào --- Trồng Cây Cho Đời Sau}
\begin{truyen}{Cụ Già Và Cây Hồ Đào --- Trồng Cây Cho Đời Sau}{Truyện Cổ La Mã --- Cicero kể lại}
\chuhoa{M}{ột} lữ khách đi ngang qua thấy một cụ già đang cần mẫn trồng cây hồ đào bên vườn.\\[4pt]
\textit{(A traveller passing by saw an old man diligently planting walnut trees by his garden.)}

Lữ khách dừng lại hỏi: ``Cụ ơi, cây này hai mươi năm nữa mới có quả; cụ còn sống đến lúc đó không?''\\[4pt]
\textit{(The traveller stopped and asked: `Old man, this tree won't bear fruit for twenty years; will you live to see it?')}

Cụ già mỉm cười đáp: ``Tôi ăn quả hồ đào cả đời mà quả đó do ai đó trước tôi trồng.''\\[4pt]
\textit{(The old man smiled and replied: `I have eaten walnuts all my life, and those trees were planted by someone before me.')}

``Đến lượt tôi, tôi trồng cho người đến sau tôi --- đó là cách thế giới này tiếp tục.''\\[4pt]
\textit{(`Now it is my turn to plant for those who come after me --- that is how this world continues.')}

Cicero kể lại câu chuyện này và nói: ``Những ai trồng cây mà họ không được hái quả là những người thật sự đang phục vụ tương lai.''\\[4pt]
\textit{(Cicero retold this story and said: `Those who plant trees under whose shade they shall never sit are truly serving the future.')}

\end{truyen}
\begin{baihoc}
  Một xã hội phát triển là xã hội mà các ông già trồng cây biết họ sẽ không ngồi dưới bóng mát của nó.\\[4pt]
  \textit{(A society grows great when old men plant trees whose shade they know they shall never sit in.)}
\end{baihoc}

\section{Andrew Carnegie Và Thư Viện Miễn Phí}
\begin{truyen}{Andrew Carnegie Và Thư Viện Miễn Phí}{Andrew Carnegie --- Hoa Kỳ, 1881}
\chuhoa{A}{ndrew} Carnegie lớn lên trong nghèo khó; trái tim ông thay đổi khi một người hàng xóm giàu có cho ông mượn sách.\\[4pt]
\textit{(Andrew Carnegie grew up in poverty; his heart changed when a wealthy neighbour allowed him to borrow books freely.)}

``Kẻ cho tôi đọc sách đó không biết rằng ông ta đang thay đổi cả cuộc đời tôi,'' ông kể lại.\\[4pt]
\textit{(`The man who gave me access to those books did not know he was changing my entire life,' he recalled.)}

Khi trở thành tỷ phú, Carnegie nhớ lại ơn đó và dùng phần lớn tài sản xây 2.509 thư viện miễn phí trên khắp thế giới.\\[4pt]
\textit{(When he became a billionaire, Carnegie remembered that kindness and used most of his wealth to build 2,509 free libraries worldwide.)}

Ông viết: ``Người giàu mà chết giàu là chết nhục --- vì họ có cơ hội làm thế giới tốt hơn mà đã không làm.''\\[4pt]
\textit{(He wrote: `The man who dies rich dies disgraced --- for he had the chance to make the world better and chose not to.')}

Hàng triệu người đọc sách trong thư viện Carnegie đã thay đổi cuộc đời --- đó là hạt giống nhân quả từ một người hàng xóm tốt bụng.\\[4pt]
\textit{(Millions who read in Carnegie's libraries changed their lives --- all seeded by one kind neighbour's generosity.)}

\end{truyen}
\begin{baihoc}
  Một hành động tốt đến đúng lúc có thể nhân lên thành hàng triệu kết quả mà bạn không bao giờ nhìn thấy được.\\[4pt]
  \textit{(One timely good deed can multiply into millions of outcomes you will never be able to see.)}
\end{baihoc}

\section{Khổng Tử Và Học Trò Nghèo Nàn Nhất}
\begin{truyen}{Khổng Tử Và Học Trò Nghèo Nàn Nhất}{Khổng Tử --- Trung Quốc, thế kỷ 5 TCN}
\chuhoa{T}{rong} số các học trò của Khổng Tử, Nhan Hồi là người nghèo nhất --- ông ở trong một căn nhà nhỏ và ăn cơm đạm bạc mỗi ngày.\\[4pt]
\textit{(Among Confucius's students, Yan Hui was the poorest --- he lived in a small hut and ate simple rice daily.)}

Nhưng mỗi khi học, Nhan Hồi luôn là người hiểu sâu nhất và thực hành nghiêm chỉnh nhất.\\[4pt]
\textit{(Yet whenever they studied, Yan Hui always understood most deeply and practised most diligently.)}

Khổng Tử hỏi ông tại sao sống khổ cực vậy mà lúc nào cũng vui; Nhan Hồi đáp:\\[4pt]
\textit{(Confucius asked why, living so humbly, he was always joyful; Yan Hui replied:)}

``Thầy dạy con rằng người quân tử tìm niềm vui trong việc làm điều phải --- con đang làm điều phải mỗi ngày.''\\[4pt]
\textit{(`Master taught me that a noble person finds joy in doing what is right --- I am doing what is right every day.')}

Khổng Tử sau đó nói với các học trò: ``Nhan Hồi là người học của ta --- và ta học được nhiều điều từ ông ấy.''\\[4pt]
\textit{(Confucius afterward told his students: `Yan Hui is my student --- and I have learned much from him.')}

\end{truyen}
\begin{baihoc}
  Người gieo hạt thiện lành mỗi ngày trong im lặng --- họ là người đang sống trong nhân quả quý nhất của cuộc đời.\\[4pt]
  \textit{(Those who sow seeds of goodness daily in silence --- they are living in the most precious cause and effect of life.)}
\end{baihoc}

\section{Matsutaro Shoriki Và Báo Yomiuri}
\begin{truyen}{Matsutaro Shoriki Và Báo Yomiuri}{Matsutaro Shoriki --- Nhật Bản, 1924}
\chuhoa{M}{atsutaro} Shoriki mua lại tờ báo Yomiuri Shimbun năm 1924 khi nó gần như phá sản hoàn toàn.\\[4pt]
\textit{(Matsutaro Shoriki acquired the Yomiuri Shimbun newspaper in 1924 when it was nearly bankrupt.)}

Ông đầu tư vào phóng viên giỏi, đầu tư vào chất lượng tin tức thay vì chạy theo lợi nhuận ngắn hạn.\\[4pt]
\textit{(He invested in skilled reporters and news quality instead of chasing short-term profit.)}

Những năm đầu lỗ nặng; nhiều người khuyên ông bán đi --- ông từ chối và tiếp tục gieo hạt chất lượng.\\[4pt]
\textit{(The early years were heavy losses; many advised him to sell --- he refused and kept sowing seeds of quality.)}

Hai mươi năm sau, Yomiuri trở thành tờ báo có số người đọc lớn nhất toàn cầu --- công của những hạt giống tốt.\\[4pt]
\textit{(Twenty years later, Yomiuri became the highest-circulation newspaper in the world --- the fruit of those good seeds.)}

Shoriki nói: ``Người gieo hạt chất lượng luôn gặt được chất lượng lại --- nhưng họ phải kiên nhẫn với mùa gặt.''\\[4pt]
\textit{(Shoriki said: `Those who sow quality always reap quality in return --- but they must be patient with the harvest.')}

\end{truyen}
\begin{baihoc}
  Gieo hạt chất lượng là hành động can đảm nhất trong nền kinh tế ngắn hạn --- vì nó đòi hỏi tin tưởng vào tương lai.\\[4pt]
  \textit{(Sowing quality is the most courageous act in a short-term economy --- for it requires faith in the future.)}
\end{baihoc}

\section{Người Thợ Rèn La Mã Và Lưỡi Kiếm}
\begin{truyen}{Người Thợ Rèn La Mã Và Lưỡi Kiếm}{Truyện Cổ La Mã --- thế kỷ 1 TCN}
\chuhoa{T}{rong} thành Rome thời cộng hoà, có một người thợ rèn nổi tiếng làm kiếm tốt nhất thành phố.\\[4pt]
\textit{(In republican Rome, there was a blacksmith famous for forging the finest swords in the city.)}

Học trò hỏi bí quyết; ông chỉ vào bể nước lạnh và nói: ``Đây là người thầy thật sự của ta.''\\[4pt]
\textit{(His apprentice asked his secret; he pointed to the cold water basin and said: `This is my true teacher.')}

Ông giải thích: ``Sắt nung nóng của ta chỉ trở nên cứng dẻo sau khi chịu đựng được nước lạnh.''\\[4pt]
\textit{(He explained: `My heated iron only becomes strong and flexible after enduring the cold water.')}

``Cuộc đời cũng vậy --- những thử thách và khó khăn không phá hủy bạn; chúng hun đúc bạn thành con người tốt hơn.''\\[4pt]
\textit{(`Life is the same --- trials and hardships do not destroy you; they forge you into a better person.')}

Mỗi lưỡi kiếm ông rèn ra mang trong mình cả triết học này --- và chúng không bao giờ gãy trong trận chiến.\\[4pt]
\textit{(Every sword he forged carried this philosophy within it --- and they never broke in battle.)}

\end{truyen}
\begin{baihoc}
  Những thử thách bạn trải qua không phải hình phạt --- chúng là quá trình rèn luyện để bạn trở nên bền vững hơn.\\[4pt]
  \textit{(The trials you go through are not punishments --- they are the forging process that makes you more resilient.)}
\end{baihoc}

\section{Martin Luther King Jr. Và Hạt Giống Bình Đẳng}
\begin{truyen}{Martin Luther King Jr. Và Hạt Giống Bình Đẳng}{Martin Luther King Jr. --- Hoa Kỳ, 1955}
\chuhoa{M}{artin} Luther King Jr. biết rằng cuộc đấu tranh dân quyền sẽ không xong trong đời ông.\\[4pt]
\textit{(Martin Luther King Jr. knew the civil rights struggle would not be completed in his lifetime.)}

Nhưng ông tiếp tục gieo hạt --- mỗi bài phát biểu, mỗi cuộc tuần hành, mỗi hành động bất tuân thủ ôn hoà.\\[4pt]
\textit{(But he kept sowing seeds --- every speech, every march, every act of peaceful non-compliance.)}

Ông nói với những người ủng hộ đang nản lòng: ``Chúng ta có thể không nhìn thấy hết Đất Hứa, nhưng chúng ta sẽ đến được.''\\[4pt]
\textit{(He told discouraged supporters: `We may not all see the Promised Land, but we, as a people, will get there.')}

Năm 1968, ông bị ám sát trước khi thấy nhiều thành quả --- nhưng hạt giống ông gieo không bao giờ dừng lớn.\\[4pt]
\textit{(In 1968 he was assassinated before seeing many fruits --- but the seeds he sowed never stopped growing.)}

Barack Obama, Tổng Thống da đen đầu tiên của Hoa Kỳ, nói trong diễn văn nhậu chức: ``Chúng ta đã đến được He.''\\[4pt]
\textit{(Barack Obama, America's first Black President, said in his inaugural address: `We have arrived.')}

\end{truyen}
\begin{baihoc}
  Hạt giống gieo vào một thế hệ thường nở hoa ở thế hệ sau --- đó không phải thất bại mà là kế hoạch của lịch sử.\\[4pt]
  \textit{(Seeds sown in one generation often bloom in the next --- that is not failure but the plan of history.)}
\end{baihoc}

\section{Người Phụ Nữ Và Bát Cơm Cứu Người}
\begin{truyen}{Người Phụ Nữ Và Bát Cơm Cứu Người}{Truyện Dân Gian Châu Á}
\chuhoa{T}{rong} một làng quê nghèo, có một người phụ nữ goá bụa mỗi ngày để dành một bát cơm ở cổng cho người qua đường đói.\\[4pt]
\textit{(In a poor village, a widowed woman every day left one bowl of rice at her gate for any passing hungry traveller.)}

Nhiều năm trôi qua; ai cũng biết về bát cơm đó nhưng không ai biết ai để lại ngoài bà.\\[4pt]
\textit{(Years passed; everyone knew of the bowl but no one knew who left it except for her.)}

Một năm mất mùa khủng khiếp ập đến; làng đói kém, bà cũng không còn gì để ăn.\\[4pt]
\textit{(One year a terrible famine struck; the village went hungry, and she herself had nothing left to eat.)}

Đúng lúc đó, cả làng kéo đến trước nhà bà --- mỗi người mang một bát thức ăn từ nhà họ để trả ơn.\\[4pt]
\textit{(At that moment, the whole village came to her door --- each person carrying a bowl of food from their own home to repay her.)}

Bà hiểu ra rằng bát cơm nhỏ bé nhiều năm của bà đã lặng lẽ gieo một mùa gặt lớn trong lòng người.\\[4pt]
\textit{(She understood that her small daily bowl over many years had quietly sown a great harvest in human hearts.)}

\end{truyen}
\begin{baihoc}
  Những hành động tốt nhỏ bé lặp đi lặp lại mỗi ngày tích luỹ thành một kho tàng lòng người lớn hơn bạn biết.\\[4pt]
  \textit{(Small good deeds repeated daily accumulate into a treasury of human goodwill greater than you know.)}
\end{baihoc}

\section{Người Đào Giếng Và Làng Khát Nước}
\begin{truyen}{Người Đào Giếng Và Làng Khát Nước}{Truyện Dân Gian Việt Nam}
\chuhoa{N}{gày} xưa có một người nông dân tốt bụng thấy cả làng mình đi xa mới có nước uống, ông quyết tâm đào giếng.\\[4pt]
\textit{(Long ago, a kind farmer seeing his whole village must travel far for water, resolved to dig a well.)}

Ông đào một mình suốt ba tháng; người trong làng ban đầu cười nhạo vì đất vùng đó rất khô cằn.\\[4pt]
\textit{(He dug alone for three months; villagers at first laughed at him for the land there was notoriously dry.)}

Đến tháng thứ tư, nước bắt đầu rỉ ra; đến tháng thứ năm, giếng đầy nước trong mát.\\[4pt]
\textit{(By the fourth month, water began seeping up; by the fifth month, the well was full of cool clear water.)}

Cả làng đổ ra uống nước mà không ai hỏi người đào là ai; nhưng người nông dân không cần ai hỏi.\\[4pt]
\textit{(The whole village came to drink and no one asked who had dug it; but the farmer needed no one to ask.)}

Ông nói với con cái: ``Giếng nước ngọt ta đào là để tặng làng --- và đó chính là điều ta sẽ không bao giờ thấy mất đi.''\\[4pt]
\textit{(He told his children: `The sweet-water well I dug is a gift to the village --- and that is something I will never see taken away.')}

\end{truyen}
\begin{baihoc}
  Hành động tốt không cần người biết tên bạn --- nó tự mang theo sức mạnh của chính mình vào thế giới.\\[4pt]
  \textit{(Good deeds do not need anyone to know your name --- they carry their own power into the world.)}
\end{baihoc}

\section{Người Đầu Bếp Nhật Và Sự Kiên Định Về Chất Lượng}
\begin{truyen}{Người Đầu Bếp Nhật Và Sự Kiên Định Về Chất Lượng}{Truyện Truyền Thống Nhật Bản}
\chuhoa{C}{ó} một người đầu bếp sushi ở Tokyo dành 10 năm chỉ học cách nấu nồi cơm --- trước khi được phép đụng vào cá.\\[4pt]
\textit{(A sushi chef in Tokyo spent 10 years learning only to cook rice --- before being permitted to touch the fish.)}

Học trò thường nản lòng và hỏi tại sao mất quá nhiều thời gian cho một việc dường như đơn giản.\\[4pt]
\textit{(Apprentices often grew discouraged and asked why so much time was spent on something seemingly simple.)}

Người thầy trả lời: ``Cơm là nền tảng; nếu nền tảng không hoàn hảo, mọi thứ trên đó đều lung lay.''\\[4pt]
\textit{(The master replied: `Rice is the foundation; if the foundation is not perfect, everything above it is unstable.')}

Học trò học cách nấu cơm mười năm rồi sau đó học mọi thứ còn lại trong ba năm --- vì nền tảng vững chắc.\\[4pt]
\textit{(The apprentice spent ten years learning rice then learned everything else in three --- because the foundation was solid.)}

Triết học ông dạy là: hãy gieo hạt kỹ năng cơ bản thật sâu trước --- mọi thứ còn lại sẽ theo sau.\\[4pt]
\textit{(The philosophy he taught was: sow the seed of basic skill deeply first --- everything else will follow.)}

\end{truyen}
\begin{baihoc}
  Gieo hạt kỹ năng nền tảng thật nghiêm túc và thật sâu --- đó là khoản đầu tư sinh lợi lâu dài nhất trong bất kỳ nghề nào.\\[4pt]
  \textit{(Sow the seeds of foundational skills seriously and deeply --- that is the longest-lasting investment in any craft.)}
\end{baihoc}
